
% =====================================================================
\section{Does 2025 look like the paper's sample?}
\label{sec:stylized}
% =====================================================================

Nothing below this point is worth reading if the measurement is wrong, and the
cheapest way to find out is to check whether a year the paper never saw
reproduces the facts it established. This section is that check. It is also the
first out-of-sample evidence on these measures.

\subsection{The headline measures}

\begin{table}[htbp]
\centering
\caption{Price clustering on the Tokyo Stock Exchange in 2024, against the published 2010--2022 benchmarks}
\label{tab:stylized}
\small
\begin{threeparttable}
\begin{tabular}{@{}lrrrrr@{}}
\toprule
Measure & 2024 mean (\%) & SE & Ohta (\%) & Difference & Stock-days \\
\midrule
$M^{0}$ & 14.07 & (0.466) & 13.5 & +0.57 & 668 \\
$M^{BLarge0}$ & 16.17 & (0.987) & 14.8 & +1.37 & 668 \\
$M^{SLarge0}$ & 14.77 & (0.559) & 14.3 & +0.47 & 668 \\
$M^{BSmall0}$ & 10.43 & (0.079) & 11.1 & -0.67 & 668 \\
$M^{SSmall0}$ & 10.30 & (0.099) & 10.9 & -0.60 & 668 \\
\bottomrule
\end{tabular}
\begin{tablenotes}[flushleft]\footnotesize
\item Volume-weighted share of trading executing at a last price digit of zero, by trade initiator and trade size. Large means a trade larger than one trading unit. Standard errors are clustered by stock and by day. The benchmark column reports Ohta's (2026) 2010--2022 means. Under a uniform digit distribution every measure would be 10\%.
\end{tablenotes}
\end{threeparttable}
\end{table}


Table~\ref{tab:stylized} puts the 2025 means beside Ohta's 2010--2022
benchmarks. Round prices take 14.22\% of all
continuous-session volume against his 13.5\%, and
14.95\% of buy-initiated large-trade volume against his 14.8\%. Under a
uniform digit distribution every one of these would be 10\%. Clustering is
alive and well on the Tokyo Stock Exchange two years past the end of his sample.

\paragraph{The ordering the mechanism predicts.}
The paper's account is specifically about \emph{large} orders taking stale
round-price limit orders, so the large-trade measure should exceed the
small-trade one. It does, by 4.18 percentage
points on the buy side, significant at the 1\% level. The small-trade measures sit
at 10.77\% and 10.95\%, barely
above the uniform benchmark. Whatever produces clustering, it is not something
that happens to trades of one unit.

\subsection{The digit distribution}

\begin{figure}[htbp]
\centering
\includegraphics[width=0.85\linewidth]{f_digits.pdf}
\caption{Volume-weighted share of trading at each last price digit, by tick
regime. The dashed line is the 10\% a uniform distribution would give.}
\label{fig:digits}
\end{figure}

Figure~\ref{fig:digits} is the whole phenomenon in one picture. Digit zero
stands well above the uniform line; digit five stands slightly above it, which is
what one expects if half-way points attract some of the same behaviour; and the
remaining seven digits sit flat. That shape is itself a diagnostic --- a
measurement error in the digit arithmetic would not produce a clean spike at zero
and a smaller one at five, it would produce something lumpy.

\subsection{Finer grids, smaller stocks}

\begin{table}[htbp]
\centering
\caption{Price clustering by tick size}
\label{tab:tick}
\small
\begin{threeparttable}
\begin{tabular}{@{}lrrrrr@{}}
\toprule
Tick (\textyen) & $M^{0}$ (\%) & $M^{BLarge0}$ (\%) & SE & Stock-days & Stocks \\
\midrule
0.1 & 18.99 & 19.05 & (0.611) & 2,095 & 30 \\
1 & 14.15 & 14.87 & (0.180) & 72,489 & 978 \\
10 & 12.45 & 13.87 & (0.486) & 2,823 & 43 \\
\bottomrule
\end{tabular}
\begin{tablenotes}[flushleft]\footnotesize
\item Ohta's sample filter admits only tick sizes that are powers of ten, since the last-digit construction is otherwise undefined. Harris (1991) predicts more clustering on finer grids.
\end{tablenotes}
\end{threeparttable}
\end{table}

\begin{table}[htbp]
\centering
\caption{Price clustering by market-capitalisation quintile}
\label{tab:size}
\small
\begin{threeparttable}
\begin{tabular}{@{}lrrrrrr@{}}
\toprule
Quintile & $M^{0}$ (\%) & $M^{BLarge0}$ (\%) & SE & $M^{BLarge0}$, \textyen 1 grid (\%) & 0.1-tick share (\%) & Stock-days \\
\midrule
Q1 & 14.45 & 14.83 & (0.304) & 14.83 & 0.0 & 15,743 \\
Q2 & 14.66 & 15.55 & (0.291) & 15.55 & 0.0 & 15,689 \\
Q3 & 14.20 & 14.77 & (0.257) & 14.77 & 0.0 & 15,683 \\
Q4 & 13.12 & 13.81 & (0.220) & 13.70 & 1.9 & 15,450 \\
Q5 & 14.70 & 15.82 & (0.250) & 15.51 & 12.2 & 14,764 \\
\bottomrule
\end{tabular}
\begin{tablenotes}[flushleft]\footnotesize
\item Quintiles are formed once per stock on its median market capitalisation over the sample, so a stock does not drift between quintiles as its price moves. Q1 is the smallest. The pooled columns mix tick regimes; the \textyen 1-grid column removes the mechanical elevation that 0.1-yen days carry, and the 0.1-tick share column shows how unevenly those days fall across quintiles.
\end{tablenotes}
\end{threeparttable}
\end{table}


\textcite{Harris1991Clustering} predicts more clustering on finer grids,
and the 0.1-yen days deliver it: 18.99\% against 14.15\% on the 1-yen grid. The mechanism is not mysterious.
At a price around 900 yen a 0.1-yen tick is about one basis point, fine enough
that order placement collapses back onto whole-yen prices --- and a whole yen is
the roundest number available. These days are kept here, where they are a
result, and excluded from the regressions, where the paper excludes them.

Table~\ref{tab:size} is the one place this sample departs from the paper. Ohta reports clustering falling with market capitalisation; here the smallest quintile shows 14.83\% and the largest 15.82\%, with no reliable gradient between them.

The likely reason is the universe rather than the market. Ohta's sample spans the whole First Section, a range wide enough for a size gradient to show. The screen used here keeps only stocks liquid enough to clear his twenty-trade and tick-size filters on more than half the sample, which removes much of the small end before the quintiles are formed; what remains is a comparison among stocks that are all reasonably active. A sample-selection explanation is not the same as a refutation, and it would be settled by running the same measures on the Standard and Growth sections, which this study does not cover.

It does not change the strategy design in Section~\ref{sec:strategy}: size is neutralised there regardless, because the concern is that a sort might load on size, not that it must.

\subsection{The impact premium}

\begin{table}[htbp]
\centering
\caption{The round-price price-impact premium in large trades}
\label{tab:dimp}
\small
\begin{threeparttable}
\begin{tabular}{@{}lrrrrrrr@{}}
\toprule
Horizon & Side & Pooled & SE & $t$ & Stock-day mean & Stock-days & Ohta \\
\midrule
1s & buy & nan & (nan) & nan & 0.254 & 1,244 & -- \\
1s & sell & nan & (nan) & nan & 0.999 & 1,479 & -- \\
60s & buy & nan & (nan) & nan & -0.091 & 1,243 & 1.26 \\
60s & sell & nan & (nan) & nan & 1.040 & 1,476 & 1.12 \\
300s & buy & nan & (nan) & nan & -0.412 & 1,239 & -- \\
300s & sell & nan & (nan) & nan & 0.805 & 1,468 & -- \\
\bottomrule
\end{tabular}
\begin{tablenotes}[flushleft]\footnotesize
\item $\Delta\mathit{Imp}$ is the volume-weighted midquote change after trades at a last price digit of zero minus the same quantity for other trades, within large trades on each side, in basis points and signed so that a positive value means the price moved the way the initiator pushed it. The pooled column removes stock and day effects and weights each stock-day by the smaller of its two cell counts; standard errors are clustered by stock and by day. Cells with fewer than five trades are dropped. The final column is Ohta's (2026) 2010--2022 mean at the one-minute horizon. $^{*}p<0.1$, $^{**}p<0.05$, $^{***}p<0.01$.
\end{tablenotes}
\end{threeparttable}
\end{table}


\begin{figure}[htbp]
\centering
\includegraphics[width=0.85\linewidth]{f_dimp.pdf}
\caption{Round-price price-impact premium in large trades, with 95\% confidence
intervals from standard errors clustered by stock and by day.}
\label{fig:dimp}
\end{figure}

This is the paper's second result and the one that makes clustering matter for
execution. A trade at a round price is executing against an order whose price has
gone stale, so the midquote should move further afterwards. At the one-minute
horizon the premium is 0.43 basis points on the buy side and
0.51 on the sell side, against Ohta's 1.26 and 1.12.

The effective-spread decomposition closes as an identity, which is a check on the
implementation rather than a finding: impact plus realized spread equals the
effective spread to three decimal places at every horizon tested.

\subsection{Standing round-price depth}

\begin{figure}[htbp]
\centering
\includegraphics[width=0.85\linewidth]{f_rdepth.pdf}
\caption{Binned means of the round-price share of visible ask depth against
measured clustering in buy-initiated large trades.}
\label{fig:rdepth}
\end{figure}

The measure the paper's data could not support says the same thing from the other
side. 13.21\% of visible ten-level ask depth sits at round prices against a
uniform benchmark of 10\%, and
14.00\% on the bid.
Figure~\ref{fig:rdepth} relates it to the executed measure: stock-days where
more of the standing book sits at round prices are stock-days where more
round-price volume executes. Where the clustering measures observe the mechanism
after the fact, this observes the inventory while it is still exposed.

\subsection{The contamination the paper warns about}

\begin{table}[htbp]
\centering
\caption{Opening-price digit, volatility, and measured clustering}
\label{tab:opening}
\small
\begin{threeparttable}
\begin{tabular}{@{}lrrrr@{}}
\toprule
Opening digit & Low volatility (\%) & High volatility (\%) & Difference & Stock-days, low vol. \\
\midrule
0 & 14.37 & 15.95 & -1.58 & 72 \\
\bottomrule
\end{tabular}
\begin{tablenotes}[flushleft]\footnotesize
\item Mean $M^{0}$ by the last digit of the opening price, split on whether the previous day's realised variance was below that day's cross-sectional median. On quiet days a stock that opens far from a round price may never reach one, which depresses the measure mechanically. Every regression in this study therefore controls for the interaction of the opening-digit dummies with the low-volatility indicator, as Ohta (2026) does.
\end{tablenotes}
\end{threeparttable}
\end{table}


Ohta's Table 2 documents a mechanical distortion: a day that opens far from a
round price and then barely moves can never print one, so its measure is low for
reasons that have nothing to do with noise traders. The distortion is present in
the sample. Opening prices cluster on their own --- digit zero accounts for
about a fifth
of openings against a tenth under uniformity --- and days opening near a round
price show visibly higher measured clustering than days opening far from one.
Every regression below therefore carries the interaction of the opening-digit
dummies with a low-volatility indicator, as the paper does.

\paragraph{Verdict.}
The sample reproduces the paper's facts: clustering well above the uniform benchmark, concentrated in large trades, stronger on finer grids, without the usual size gradient, carrying a positive round-price impact premium. The measurement is sound and the rest of the study can be built on it.
